% sections/term_margin.tex: Section 5.4: the term margin as three channels, and the deficiency and recourse contrasts that
% identify the penalty channel. Results: identification_profile_by_decile.csv (I6c), identification_horizon_profile.csv (I6b),
% identification_deficiency_powered.csv (I3b), identification_within_person.csv (I6).
\subsection{The Term Margin and the Cost of Default}\label{sec:termmargin}\label{sec:termstructure}

The theory's second margin is the response of default to payments added at the far end of the contract. In the credit-bureau data that margin is the coefficient on the contract's length at a fixed monthly payment, and Proposition \ref{prop:termmargin} says it is three things: the weight on the added payments, minus the cost of the larger balance, plus the defaults of the added months when default is measured over the loan's life. Table \ref{tab:term_by_score} reports it by credit-score decile and by the horizon over which default is measured, within lender-year and zip3-by-year cells, and in the last row within borrower. The table has three regularities.

First, at a fixed payment a longer contract lowers early default and raises lifetime default. With the borrower held fixed, on every multi-loan borrower in the ten-percent archive, the elasticity of default to the term is $-0.31$ within twelve months, $+0.30$ at 24, and $+0.95$ over the life of the loan, while the payment elasticity is $0.44$, $0.53$ and $0.46$ (Table \ref{tab:main_margins}, lower panel). Absorbing lender-year in place of origination-year cells on the same loans, which holds the lender's menu fixed as well as the borrower, gives the last row of Table \ref{tab:term_by_score}: a term elasticity of $-0.44$ within six months and $-0.19$ within twelve, $0.00$ ($0.01$) at twenty-four and $+0.42$ over the life, with payment elasticities of $0.51$ and $0.57$ (Appendix \ref{appendix:ivvalidation}); the paper's estimates use the origination-year cells, which compare a borrower's contracts across lenders at the same date. Second, the early effect is largest for the riskiest borrowers and reverses sign near the middle of the score distribution: within twenty-four months the term elasticity is $-0.53$ in the bottom decile and $+0.57$ in the sixth, with each intermediate decile ordered between them. Third, on the lifetime outcome the elasticity rises with the score from $0.12$ to above one.

The lifetime pattern is exposure: a longer contract has more months in which a default can occur, so the coefficient on the term in any outcome measured over the contract's life contains the additional hazards of the added months, an effect with no content about preferences. The early pattern is what remains once exposure is removed by fixing the window, and it has the wrong sign for a model in which the cost of default is a fixed penalty. It has the right sign for a model in which the cost of default rises with the balance owed. A longer contract at a fixed payment is a larger balance at every date; if the lender can pursue the shortfall after taking the collateral, or the credit record reports the size of the charge-off, defaulting early costs more, and default falls with the term. Appendix \ref{appendix:balancepenalty} derives the term margin under a penalty that scales with the balance: the response to a payment at date $s$ is the preference weight $D(s-1)Q_s\lambda$ minus a penalty weight $\kappa_B(1+r)^{-(s-1)}$, the first discounted at the borrower's rate and the second at the contract's. The penalty weight is discounted only at the contract rate, so the penalty channel can dominate at short horizons whatever the borrower's rate, and the negative early term effect says that it does; the ordering reverses as the loan amortizes. That produces the horizon profile. It produces the score gradient because the borrowers whose defaults are early are observed at the horizons where the penalty dominates.

\paragraph{The deficiency and recourse contrasts.} The prediction that separates this account from selection is that the early effect should be weaker where the balance cannot be pursued, and the law of the states supplies the variation. Sixteen states bar a deficiency judgment on an auto loan below a threshold balance; the design compares loans below the threshold in those states with loans of the same size in states without a bar, within zip3-by-origination-year cells (Table \ref{tab:term_deficiency}). On auto loans below \$8,100 the term elasticity of default within twelve months is $-1.01$ and is $0.36$ (standard error $0.13$) less negative in the states with a bar; at 36 months the difference is $1.02$ ($0.27$). On first mortgages, using the eleven states whose law bars recourse to the borrower's other assets, the difference is $3.9$ ($1.5$) at twelve months and $2.2$ ($1.0$) at 24. Every cell has the predicted sign. Selection on the borrower's stable creditworthiness is separately excluded: within a lender's book higher scores go with shorter terms at a given payment, adding score, income and age to the regression makes the term coefficient more negative rather than less, and the sign is unchanged with borrower fixed effects. The contrast identifies the part of the penalty channel that the deficiency judgment accounts for; the credit record's report of the charge-off is the other part, and Appendix \ref{appendix:ladder} states how the two are combined when the across-borrower term margin is adjusted.

The consequence for measurement is direct. The term margin from installment contracts is the difference of two channels plus exposure, and the preference channel is the smallest of the three at the horizons the data cover. A ratio of the term response to the payment response identifies the discount function only with the penalty scaling and the exposure removed, and a measured ratio above one is outside the preference channel at every rate. The term margin is therefore used in two ways: within borrower at twenty-four months, where exposure is removed by the window and the penalty channel has decayed with the balance, as the far-payment response the rate rests on (Section \ref{sec:rate_people}); and across borrowers, without a model, as evidence on how the cost of default scales with what is owed, which is what Table \ref{tab:term_deficiency} reports. Default is a strategic decision, responsive to the balance at stake and to whether the lender can collect it. The estimator extends to revolving credit, where no contract fixes the payment schedule and the two margins are defined by the account's position relative to the issuer's minimum-payment formula; Appendix \ref{appendix:cards} states that design and its feasibility in the archives.

\begin{table}[ht]
\centering
\caption{The term margin by credit-score decile and horizon (auto loans)}
\label{tab:term_by_score}
\resizebox{\textwidth}{!}{\begin{tabular}{lcccccccc}
\toprule
 & \multicolumn{5}{c}{Term elasticity, default within} & \multicolumn{3}{c}{Payment elasticity} \\
\cmidrule(lr){2-6} \cmidrule(lr){7-9}
Score decile & 6 mo & 12 mo & 24 mo & 36 mo & life & 12 mo & 24 mo & life \\
\midrule
1 & -1.07 (0.06) & -0.75 (0.04) & -0.53 (0.04) & -0.45 (0.04) & +0.12 (0.01) & +0.26 (0.02) & +0.29 (0.01) & +0.19 (0.01) \\
2 & -0.90 (0.06) & -0.61 (0.04) & -0.40 (0.04) & -0.29 (0.04) & +0.20 (0.02) & +0.28 (0.02) & +0.33 (0.02) & +0.24 (0.01) \\
3 & -0.56 (0.06) & -0.36 (0.04) & -0.14 (0.03) & -0.05 (0.04) & +0.34 (0.02) & +0.35 (0.02) & +0.40 (0.02) & +0.30 (0.01) \\
4 & -0.33 (0.05) & -0.15 (0.04) & +0.08 (0.04) & +0.17 (0.05) & +0.48 (0.03) & +0.39 (0.02) & +0.45 (0.02) & +0.36 (0.02) \\
5 & +0.01 (0.06) & +0.18 (0.04) & +0.36 (0.04) & +0.46 (0.05) & +0.70 (0.03) & +0.48 (0.03) & +0.51 (0.02) & +0.39 (0.02) \\
6 & +0.23 (0.09) & +0.44 (0.05) & +0.57 (0.05) & +0.67 (0.05) & +0.92 (0.04) & +0.48 (0.03) & +0.49 (0.03) & +0.39 (0.02) \\
7 & +0.43 (0.12) & +0.71 (0.07) & +0.93 (0.05) & +1.06 (0.06) & +1.20 (0.04) & +0.44 (0.04) & +0.43 (0.03) & +0.34 (0.02) \\
8 & +0.53 (0.15) & +0.87 (0.09) & +1.03 (0.07) & +1.11 (0.08) & +1.35 (0.05) & +0.27 (0.05) & +0.26 (0.03) & +0.18 (0.02) \\
\midrule
all, within borrower & -0.44 (0.02) & -0.19 (0.01) & +0.00 (0.00) & +0.07 (0.00) & +0.42 (0.00) & +0.51 (0.00) & +0.57 (0.00) & +0.47 (0.00) \\
\bottomrule
\end{tabular}}
\vspace{0.5em}
\begin{minipage}{0.95\textwidth}
\small \textit{Notes:} Elasticities of the default rate with respect to the contract term and the monthly payment from a linear probability model of default within the stated window on $\log m$ and $\log T$, within lender-year and zip3-by-origination-year cells, on the covariate-matched ten-percent auto extract (72,906 lifetime defaults), the sample on which the deficiency contrast runs; the panel's payment-elasticity profile by decile is in Section \ref{sec:heterogeneity}. Standard errors clustered by lender-year. The last row is the within-borrower specification with borrower and lender-year fixed effects on the multi-loan borrowers of the ten-percent archive. Deciles above the eighth have too few early defaults to estimate.
\end{minipage}
% replication: Code/identification/I6c_profile_by_decile_full.R, Code/identification/I6b_horizon_profile_by_group.R, Code/exhibits/make_term_margin_tables.py
\end{table}

\begin{table}[ht]
\centering
\caption{The early term effect where the balance cannot be pursued}
\label{tab:term_deficiency}
\resizebox{\textwidth}{!}{\begin{tabular}{llcccrr}
\toprule
Sample & Outcome & Term elasticity & $\times$ balance not pursuable & Payment elasticity & Defaults & Share no-pursuit, \% \\
\midrule
Autos, loans under \$8,100 (small-balance bars) & 12 mo & -1.01 (0.12) & +0.36 (0.13) & +1.15 & 63,949 & 6.0 \\
 & 24 mo & -0.10 (0.06) & +0.38 (0.35) & +1.82 & 95,495 & 5.5 \\
 & 36 mo & +0.96 (0.09) & +1.02 (0.27) & +2.08 & 73,462 & 4.8 \\
 & life & +0.22 (0.04) & +0.20 (0.12) & +1.24 & 161,230 & 6.0 \\
Autos, all loans (triple difference) & 12 mo & -0.12 (0.09) & +0.51 (0.22) & +0.22 & 405,464 & 0.6 \\
 & 24 mo & +0.33 (0.06) & +0.22 (0.38) & +0.34 & 955,840 & 0.5 \\
 & 36 mo & +0.64 (0.05) & +0.63 (0.19) & +0.35 & 1,266,466 & 0.3 \\
 & life & +0.94 (0.05) & +0.11 (0.10) & +0.27 & 2,124,154 & 0.6 \\
Mortgages (non-recourse states) & 12 mo & -8.23 (1.14) & +3.87 (1.52) & -4.59 & 3,472 & 30.3 \\
 & 24 mo & -5.33 (0.78) & +2.17 (1.01) & -3.40 & 6,831 & 30.4 \\
 & 36 mo & -3.05 (0.49) & +0.77 (0.67) & -2.35 & 10,055 & 30.3 \\
 & life & -2.72 (0.36) & +0.73 (0.51) & -2.31 & 30,715 & 30.2 \\
\bottomrule
\end{tabular}}
\vspace{0.5em}
\begin{minipage}{0.95\textwidth}
\small \textit{Notes:} Linear probability models of default within the stated window on $\log m$, $\log T$, and their interactions with an indicator that the unpaid balance cannot be pursued after the collateral is taken, within zip3-by-origination-year cells; standard errors clustered by state. Autos: loans below the state's small-balance deficiency threshold, compared with loans of the same size in states without one. Mortgages: the eleven non-recourse states of \citet{ghent_recourse_2011}.
\end{minipage}
% replication: Code/identification/I3b_deficiency_powered.R, Code/exhibits/make_term_margin_tables.py; external/state_deficiency_rules.csv
\end{table}
